\documentclass[12pt]{article}
\usepackage[paperwidth=612bp,paperheight=792bp,margin=0bp,noheadfoot]{geometry}
\usepackage{fontspec}
\usepackage{xeCJK}
\usepackage{tikz}
\pagestyle{empty}
\setlength{\parindent}{0pt}
\setlength{\parskip}{0pt}
\IfFontExistsTF{TeX Gyre Termes}{\setmainfont{TeX Gyre Termes}}{\setmainfont{Times New Roman}}
\IfFontExistsTF{Noto Serif CJK SC}{\setCJKmainfont{Noto Serif CJK SC}}{\setCJKmainfont{Songti SC}}
\newcommand{\QuestionTitle}{人机实验1 H-C2：项羽本纪}
\newcommand{\DrawQuestionTitle}{%
\node[anchor=base west,inner sep=0pt,outer sep=0pt] at ([xshift=72bp,yshift=-34bp]current page.north west) {{\fontsize{14bp}{14bp}\selectfont\bfseries \QuestionTitle}};%
\node[anchor=base west,inner sep=0pt,outer sep=0pt] at ([xshift=72bp,yshift=-62bp]current page.north west) {{\fontsize{12bp}{12bp}\selectfont 用时：\underline{\hspace{80bp}}}};%
}
\begin{document}
% Auto-generated subset for 项羽
\thispagestyle{empty}
\begin{tikzpicture}[remember picture,overlay]
\DrawQuestionTitle
\node[anchor=base west,inner sep=0pt,outer sep=0pt] at ([xshift=72bp,yshift=-84bp]current page.north west) {{\fontsize{12bp}{12bp}\selectfont 【7.25.1】 汉元年四月，在项王硅麾之下诸侯罢兵散归，各自回到封国。}};
\node[anchor=base west,inner sep=0pt,outer sep=0pt] at ([xshift=72bp,yshift=-100bp]current page.north west) {{\fontsize{12bp}{12bp}\selectfont 【7.25.2】 项王也出关回到封国，派人迁徙义帝，说：“古代做帝王的拥有千里见方的土}};
\node[anchor=base west,inner sep=0pt,outer sep=0pt] at ([xshift=72bp,yshift=-116bp]current page.north west) {{\fontsize{12bp}{12bp}\selectfont 地，必须住在上游。”于是就派遣使者把义帝迁往长沙郴县。}};
\node[anchor=base west,inner sep=0pt,outer sep=0pt] at ([xshift=72bp,yshift=-132bp]current page.north west) {{\fontsize{12bp}{12bp}\selectfont 【7.25.3】 项王催促义帝快些动身，义帝群臣渐渐背叛了他，项王就暗中命令衡山王、临}};
\node[anchor=base west,inner sep=0pt,outer sep=0pt] at ([xshift=72bp,yshift=-148bp]current page.north west) {{\fontsize{12bp}{12bp}\selectfont 江王把义帝击杀在江中。}};
\node[anchor=base west,inner sep=0pt,outer sep=0pt] at ([xshift=72bp,yshift=-164bp]current page.north west) {{\fontsize{12bp}{12bp}\selectfont 【7.25.4】 韩王成没有军功，项王不让他就国，一起到了彭城，废去王号，改封为侯，不}};
\node[anchor=base west,inner sep=0pt,outer sep=0pt] at ([xshift=72bp,yshift=-180bp]current page.north west) {{\fontsize{12bp}{12bp}\selectfont 久又杀死了。}};
\node[anchor=base west,inner sep=0pt,outer sep=0pt] at ([xshift=72bp,yshift=-196bp]current page.north west) {{\fontsize{12bp}{12bp}\selectfont 【7.25.5】 臧荼到了封国，就驱逐韩广去辽东，韩广不服从，臧茶在无终击杀了韩广，兼}};
\node[anchor=base west,inner sep=0pt,outer sep=0pt] at ([xshift=72bp,yshift=-212bp]current page.north west) {{\fontsize{12bp}{12bp}\selectfont 并了他的封地。}};
\node[anchor=base west,inner sep=0pt,outer sep=0pt] at ([xshift=72bp,yshift=-228bp]current page.north west) {{\fontsize{12bp}{12bp}\selectfont 【7.26.1】 田荣听说项羽把齐王市徒封胶东，而立齐将田都为齐王，十分气愤，不愿让齐}};
\node[anchor=base west,inner sep=0pt,outer sep=0pt] at ([xshift=72bp,yshift=-244bp]current page.north west) {{\fontsize{12bp}{12bp}\selectfont 王去胶东，就据齐反叛，迎击田都。}};
\node[anchor=base west,inner sep=0pt,outer sep=0pt] at ([xshift=72bp,yshift=-260bp]current page.north west) {{\fontsize{12bp}{12bp}\selectfont 【7.26.2】 田都逃往楚国。}};
\node[anchor=base west,inner sep=0pt,outer sep=0pt] at ([xshift=72bp,yshift=-276bp]current page.north west) {{\fontsize{12bp}{12bp}\selectfont 【7.26.3】 齐工市害怕项王，就潜往胶东就国。}};
\node[anchor=base west,inner sep=0pt,outer sep=0pt] at ([xshift=72bp,yshift=-292bp]current page.north west) {{\fontsize{12bp}{12bp}\selectfont 【7.26.4】 田荣大为生气，派兵追击，在即墨杀死了他。}};
\node[anchor=base west,inner sep=0pt,outer sep=0pt] at ([xshift=72bp,yshift=-308bp]current page.north west) {{\fontsize{12bp}{12bp}\selectfont 【7.26.5】 田荣便自立为齐王，向西进兵，击杀了济北王田安，兼并了三齐。}};
\node[anchor=base west,inner sep=0pt,outer sep=0pt] at ([xshift=72bp,yshift=-324bp]current page.north west) {{\fontsize{12bp}{12bp}\selectfont 【7.26.6】 田荣把将军印授予彭越，让他在梁地反楚。}};
\node[anchor=base west,inner sep=0pt,outer sep=0pt] at ([xshift=72bp,yshift=-340bp]current page.north west) {{\fontsize{12bp}{12bp}\selectfont 【7.26.7】 陈余秘密派遣张同、夏说劝告齐王田荣说：“项羽为天下的主宰，（分封侯王}};
\node[anchor=base west,inner sep=0pt,outer sep=0pt] at ([xshift=72bp,yshift=-356bp]current page.north west) {{\fontsize{12bp}{12bp}\selectfont ）不公平。如今把原来的诸侯王都封在坏地方称王，而他的群臣诸将都封在好地方称王。}};
\node[anchor=base west,inner sep=0pt,outer sep=0pt] at ([xshift=72bp,yshift=-372bp]current page.north west) {{\fontsize{12bp}{12bp}\selectfont （因为要）赶走原来的诸侯王，赵王就（只好）到北方居住代地，我以为这样是不能答应}};
\node[anchor=base west,inner sep=0pt,outer sep=0pt] at ([xshift=72bp,yshift=-388bp]current page.north west) {{\fontsize{12bp}{12bp}\selectfont 的。听说大王已经起兵，而且不接受不道义的命令，希望大王援助我一些兵马，允许我用}};
\node[anchor=base west,inner sep=0pt,outer sep=0pt] at ([xshift=72bp,yshift=-404bp]current page.north west) {{\fontsize{12bp}{12bp}\selectfont 以攻打常山，恢复赵王的地位，愿把赵国作为齐国的屏障。”}};
\node[anchor=base west,inner sep=0pt,outer sep=0pt] at ([xshift=72bp,yshift=-420bp]current page.north west) {{\fontsize{12bp}{12bp}\selectfont 【7.26.8】 齐王答应了，就遣兵赴赵。}};
\node[anchor=base west,inner sep=0pt,outer sep=0pt] at ([xshift=72bp,yshift=-436bp]current page.north west) {{\fontsize{12bp}{12bp}\selectfont 【7.26.9】 陈余调动了三县的全部士卒，与齐军并力攻打常山，打垮了常山的军队。}};
\node[anchor=base west,inner sep=0pt,outer sep=0pt] at ([xshift=72bp,yshift=-452bp]current page.north west) {{\fontsize{12bp}{12bp}\selectfont 【7.26.10】 张耳逃走归服了汉王。}};
\node[anchor=base west,inner sep=0pt,outer sep=0pt] at ([xshift=72bp,yshift=-468bp]current page.north west) {{\fontsize{12bp}{12bp}\selectfont 【7.26.11】 陈余去代地迎接原来的赵王歇返归赵地。}};
\node[anchor=base west,inner sep=0pt,outer sep=0pt] at ([xshift=72bp,yshift=-484bp]current page.north west) {{\fontsize{12bp}{12bp}\selectfont 【7.26.12】 赵王就立陈余为代王。}};
\node[anchor=base west,inner sep=0pt,outer sep=0pt] at ([xshift=72bp,yshift=-500bp]current page.north west) {{\fontsize{12bp}{12bp}\selectfont 【7.27.1】 这时，汉王回军平定了三秦。}};
\node[anchor=base west,inner sep=0pt,outer sep=0pt] at ([xshift=72bp,yshift=-516bp]current page.north west) {{\fontsize{12bp}{12bp}\selectfont 【7.27.2】 项羽听说汉王已经兼并了关中，将要东进，齐、赵又反叛了他，非常愤怒。}};
\node[anchor=base west,inner sep=0pt,outer sep=0pt] at ([xshift=72bp,yshift=-532bp]current page.north west) {{\fontsize{12bp}{12bp}\selectfont 【7.27.3】 就以从前的吴令郑昌为韩王，来阻挡汉军。}};
\node[anchor=base west,inner sep=0pt,outer sep=0pt] at ([xshift=72bp,yshift=-548bp]current page.north west) {{\fontsize{12bp}{12bp}\selectfont 【7.27.4】 命令萧公角等人攻击彭越。}};
\node[anchor=base west,inner sep=0pt,outer sep=0pt] at ([xshift=72bp,yshift=-564bp]current page.north west) {{\fontsize{12bp}{12bp}\selectfont 【7.27.5】 彭越打败了萧公角等人。}};
\node[anchor=base west,inner sep=0pt,outer sep=0pt] at ([xshift=72bp,yshift=-580bp]current page.north west) {{\fontsize{12bp}{12bp}\selectfont 【7.27.6】 汉王派张良巡行招抚韩地，张良就给项王写信说：“汉王（没有如约称王关中}};
\node[anchor=base west,inner sep=0pt,outer sep=0pt] at ([xshift=72bp,yshift=-596bp]current page.north west) {{\fontsize{12bp}{12bp}\selectfont ），有失职守，打算取得关中，实现了原来的约定就停止进军，不敢继续东进。”}};
\node[anchor=base west,inner sep=0pt,outer sep=0pt] at ([xshift=72bp,yshift=-612bp]current page.north west) {{\fontsize{12bp}{12bp}\selectfont 【7.27.7】 张良又把齐、梁的反叛文告送给项王，说：“齐想和赵并力灭楚。”楚军因此}};
\node[anchor=base west,inner sep=0pt,outer sep=0pt] at ([xshift=72bp,yshift=-628bp]current page.north west) {{\fontsize{12bp}{12bp}\selectfont 无意西进，而向北攻打齐国。}};
\node[anchor=base west,inner sep=0pt,outer sep=0pt] at ([xshift=72bp,yshift=-644bp]current page.north west) {{\fontsize{12bp}{12bp}\selectfont 【7.27.8】 项王向九江王黥布征调兵力。}};
\node[anchor=base west,inner sep=0pt,outer sep=0pt] at ([xshift=72bp,yshift=-660bp]current page.north west) {{\fontsize{12bp}{12bp}\selectfont 【7.27.9】 黥布称病不往，派将领率兵几千人前去。}};
\node[anchor=base west,inner sep=0pt,outer sep=0pt] at ([xshift=72bp,yshift=-676bp]current page.north west) {{\fontsize{12bp}{12bp}\selectfont 【7.27.10】 项王从此怨恨黥布。}};
\node[anchor=base west,inner sep=0pt,outer sep=0pt] at ([xshift=72bp,yshift=-692bp]current page.north west) {{\fontsize{12bp}{12bp}\selectfont 【7.28.1】 汉二年冬，项羽北上到达城阳，田荣也率军到此与项羽会战。}};
\node[anchor=base west,inner sep=0pt,outer sep=0pt] at ([xshift=72bp,yshift=-708bp]current page.north west) {{\fontsize{12bp}{12bp}\selectfont 【7.28.2】 田荣兵败，逃到平原，平原百姓杀死了他。}};
\node[anchor=base west,inner sep=0pt,outer sep=0pt] at ([xshift=72bp,yshift=-724bp]current page.north west) {{\fontsize{12bp}{12bp}\selectfont 【7.28.3】 楚军北进，烧毁齐国房屋，夷平齐国城郭，坑杀田荣降卒，掳掠老弱妇女。}};
\end{tikzpicture}
\null
\newpage
\thispagestyle{empty}
\begin{tikzpicture}[remember picture,overlay]
\DrawQuestionTitle
\node[anchor=base west,inner sep=0pt,outer sep=0pt] at ([xshift=72bp,yshift=-84bp]current page.north west) {{\fontsize{12bp}{12bp}\selectfont 【7.32.1】 汉三年，项王屡次侵夺汉军的甬道，汉王粮食缺乏，恐慌起来，请求讲和，划}};
\node[anchor=base west,inner sep=0pt,outer sep=0pt] at ([xshift=72bp,yshift=-100bp]current page.north west) {{\fontsize{12bp}{12bp}\selectfont 分荥阳以西归汉。}};
\node[anchor=base west,inner sep=0pt,outer sep=0pt] at ([xshift=72bp,yshift=-116bp]current page.north west) {{\fontsize{12bp}{12bp}\selectfont 【7.32.2】 项王想要答应他。}};
\node[anchor=base west,inner sep=0pt,outer sep=0pt] at ([xshift=72bp,yshift=-132bp]current page.north west) {{\fontsize{12bp}{12bp}\selectfont 【7.32.3】 历阳侯范增说：“汉军容易对付，现在放掉他们，不予以消灭，以后一定要懊}};
\node[anchor=base west,inner sep=0pt,outer sep=0pt] at ([xshift=72bp,yshift=-148bp]current page.north west) {{\fontsize{12bp}{12bp}\selectfont 悔。”}};
\node[anchor=base west,inner sep=0pt,outer sep=0pt] at ([xshift=72bp,yshift=-164bp]current page.north west) {{\fontsize{12bp}{12bp}\selectfont 【7.32.4】 项王就和范增加紧围攻荥阳。}};
\node[anchor=base west,inner sep=0pt,outer sep=0pt] at ([xshift=72bp,yshift=-180bp]current page.north west) {{\fontsize{12bp}{12bp}\selectfont 【7.35.2】 汉四年，项王围攻成皋。}};
\end{tikzpicture}
\null
\end{document}

